\documentclass[a4paper]{article}

\usepackage{silence}
\WarningsOff[everypage]

\usepackage[pages=all, color=black, position={current page.south}, placement=bottom, scale=1, opacity=1, vshift=5mm]{background}
\SetBgContents{
	\tt This work is shared under a \href{https://creativecommons.org/licenses/by-sa/4.0/}{CC BY-SA 4.0 license} unless otherwise noted
}      % copyright

\usepackage[margin=1in]{geometry} % full-width

\usepackage[T1]{fontenc}
\usepackage{parskip}
\parskip=10pt
\usepackage[english]{isodate}

% AMS Packages
\usepackage{amsmath}
\usepackage{amsthm}
\usepackage{amssymb}

% Unicode
\usepackage[utf8]{inputenc}
\usepackage{hyperref}
\hypersetup{
	unicode,
	colorlinks,
	breaklinks,
	urlcolor=cyan, 
	linkcolor=blue, 
	pdfauthor={Author One, Author Two, Author Three},
	pdftitle={A simple article template},
	pdfsubject={A simple article template},
	pdfkeywords={article, template, simple},
	pdfproducer={LaTeX},
	pdfcreator={pdflatex}
}

% Natbib
\usepackage[sort&compress,numbers,square]{natbib}
\bibliographystyle{mplainnat}

% Theorem, Lemma, etc
\theoremstyle{plain}
\newtheorem{theorem}{Theorem}[section]
\newtheorem{corollary}[theorem]{Corollary}
\newtheorem{lemma}[theorem]{Lemma}
\newtheorem{claim}{Claim}[theorem]
\newtheorem{axiom}[theorem]{Axiom}
\newtheorem{conjecture}[theorem]{Conjecture}
\newtheorem{fact}[theorem]{Fact}
\newtheorem{hypothesis}[theorem]{Hypothesis}
\newtheorem{assumption}[theorem]{Assumption}
\newtheorem{proposition}[theorem]{Proposition}
\newtheorem{criterion}[theorem]{Criterion}
\theoremstyle{definition}
\newtheorem{definition}[theorem]{Definition}
\newtheorem{example}[theorem]{Example}
\newtheorem{remark}[theorem]{Remark}
\newtheorem{problem}[theorem]{Problem}
\newtheorem{principle}[theorem]{Principle}

\usepackage{graphicx, color}
\graphicspath{{fig/}}

%\usepackage[linesnumbered,ruled,vlined,commentsnumbered]{algorithm2e} % use algorithm2e for typesetting algorithms
\usepackage{algorithm, algpseudocode} % use algorithm and algorithmicx for typesetting algorithms
\usepackage{mathrsfs} % for \mathscr command

\usepackage{lipsum}

% Author info
\title{Linear Algebra Applications - Discrete Linear Dynamical Systems}
\author{Alexander Xie$^1$\thanks{With support from Pau Battle}\\\\
	\href{mailto:alex.xie1@student.unsw.edu.au}{alex.xie1@student.unsw.edu.au}}

\date{\printdayoff\today}

\begin{document}

\maketitle

\begin{abstract}
	Dynamical systems are processes that evolve in time. Through forms of mathematical modelling, the study of dynamical
	systems is beneficial to comprehend a variety of natural phenomena. In particular, discrete linear dynamical systems
	are useful to model a scenario as a series of state snapshots.
	In this paper, we focus on discrete linear dynamical systems in the form of first-order homogenous matrix recurrence
	relation. We then illustrate the application of these systems through a real-world case study by modelling the
	spread of a virus.

	% \noindent\textbf{Keywords:} article, template, simple
\end{abstract}

\tableofcontents

\pagebreak

\section{Discrete Linear Dynamical Systems}

\subsection{A matrix recurrence relation}

For a square matrix $A$ and an initial state vector $x_0$, define the first-order homogenous recurrence:
\[ x_t=Ax_{t-1},\quad \text{for } t\geq 1. \]
Iterating yields a sequence of vectors $x_0, Ax_0, A^2x_0,\ldots, A^kx_0,\ldots$ which represent the evolution of the
system starting from $x_0$. This sequence is called the \textbf{\textit{orbit}} of the
system starting from $x_0$.

An important question to raise about this dynamical system is: what is its long-term behaviour? For example, if our
system models a particular population of rabbits, it would be worthwhile to observe the population a few months ahead.
What happens as $t\to\infty$? Does the system have limiting behaviours?

To compute the elements in our orbit, we require an efficient method to compute $A^n$ for large $n$, as repeated
matrix multiplication is cumbersome.

\subsection{Diagonalisation}

Recall that it is easy to multiply diagonal matrices (square matrices in which the entries outside the main diagonal
are all zero):

\[
	\begin{pmatrix}
		a_{11} &        &        &        \\
		       & a_{22} &        &        \\
		       &        & \ddots &        \\
		       &        &        & a_{nn}
	\end{pmatrix}
	\begin{pmatrix}
		b_{11} &        &        &        \\
		       & b_{22} &        &        \\
		       &        & \ddots &        \\
		       &        &        & b_{nn}
	\end{pmatrix}
	=
	\begin{pmatrix}
		a_{11}b_{11} &              &        &              \\
		             & a_{22}b_{22} &        &              \\
		             &              & \ddots &              \\
		             &              &        & a_{nn}b_{nn}
	\end{pmatrix}.
\]

This simplicity of multiplying diagonal matrices prompts us to investigate if $A$ is similar to a diagonal matrix; does
there exist an invertible matrix $P$ and a diagonal matrix $D$ such that $A$ can be written as $PDP^{-1}$? If so, then
we say that $A$ is \textbf{\textit{diagonalisable}}. This condition is equivalent to $AP=PD$. Observe that it is easier
to computer powers of a diagonalisable matrix:
\begin{align*}
	A^n & = (PDP^{-1})(PDP^{-1})\cdots(PDP^{-1})                                                             \\
	    & = PD\underbrace{(P^{-1}P)}_{I}D\underbrace{(P^{-1}P)}_{I}D\cdots \underbrace{(P^{-1}P)}_{I}DP^{-1} \\
	    & = PD^n P^{-1}.
\end{align*}

\begin{theorem}
	\label{thrm:1}
	% The matrix $A\in\mathbb{R}^{n\times n}$ is diagonalisable if and only if $A$ has $n$ linearly independent eigenvectors.
	A square matrix is diagonalisable if and only if it has $n$ linearly independent eigenvectors.
\end{theorem}

\begin{proof}
	$(\implies)$ Suppose that $A\in\mathbb{R}^{n\times n}$ is diagonalisable. Then there exists an invertible matrix
	$P\in\mathbb{R}^{n\times n}$ and a diagonal matrix $D\in\mathbb{R}^{n\times n}$ such that $AP=PD$.

	Let $P=\begin{pmatrix}v_1 & v_2 & \cdots & v_n \end{pmatrix}$ where $v_i$ are linearly independent columns, and let
	$D=\operatorname{diag}(\lambda_1,\lambda_2,\ldots,\lambda_n)$.

	Writing $AP=PD$ as columns, we have $\begin{pmatrix}Av_1 & Av_2 & \cdots & Av_n\end{pmatrix}=
		\begin{pmatrix}\lambda_1v_1 & \lambda_2v_2 & \cdots & \lambda_nv_n\end{pmatrix}$.

	Thus, $Av_i=\lambda_iv_i$ for each $i=1,\ldots n$. Hence, $v_i$ is an eigenvector associated to $\lambda_i$.

	$(\impliedby)$ Conversely, suppose that $A\in\mathbb{R}^{n\times n}$ has $n$ linearly independent eigenvectors
	$v_1,v_2,\ldots,v_n$ associated to the eigenvalues $\lambda_1,\lambda_2,\ldots,\lambda_n$ respectively. Then for each
	$i=1,\ldots,n$ we have $Av_i=\lambda_iv_i$.

	Let $P=\begin{pmatrix}v_1 & v_2 & \cdots & v_n\end{pmatrix}$ where $v_i$ are the columns. Note that $P$ is invertible
	as the columns are linearly independent. Also, let $D=\operatorname{diag}(\lambda_1,\lambda_2,\cdots\lambda_n)$.

	Now the equations $Av_i=\lambda_iv_i$ for each $i=1,\ldots,n$ can be compactly expressed as $AP=PD$, which implies
	$A=PDP^{-1}$. Hence, $A$ is diagonalisable.
\end{proof}

\begin{corollary}
	\label{cor:2}
	An eigen-formula for diagonalisable matrices.
\end{corollary}

In \ref{thrm:1}, observe that $P$ and $D$ are constructed by the eigenvectors and eigenvalues of $A$. So if
$A\in\mathbb{R}^{n\times n}$ has $n$ linearly independent eigenvectors $v_1,v_2,\ldots,v_n$ associated to the eigenvalues
$\lambda_1,\lambda_2,\ldots,\lambda_n$, then

\[
	A=\begin{pmatrix}v_1 & v_2 & \cdots & v_n\end{pmatrix}
	\begin{pmatrix}
		\lambda_1 &           &        &           \\
		          & \lambda_2 &        &           \\
		          &           & \ddots &           \\
		          &           &        & \lambda_n
	\end{pmatrix}
	\begin{pmatrix}v_1 & v_2 & \cdots & v_n\end{pmatrix}^{-1}.
\]

\begin{example}
	Diagonalise $A=\begin{pmatrix}2&0&0\\-2&1&1\\1&0&-1\end{pmatrix}$, if possible.

	We will first find the eigenvalues of $A$ by solving $\det(\lambda I - A)=0$:
	\begin{align*}
		0 & =\det\begin{pmatrix}\lambda-2&0&0\\2&\lambda-1&-1\\-1&0&\lambda+1\end{pmatrix} \\
		  & =(\lambda-2)\det\begin{pmatrix}\lambda-1&-1\\0&\lambda+1\end{pmatrix}          \\
		  & =(\lambda-2)(\lambda-1)(\lambda+1).
	\end{align*}
	Thus the eigenvalues of $A$ are $\lambda_1=1,\lambda_2=2,\lambda_3=-1$.

	To find the eigenvectors of $A$, we solve $(\lambda_i I - A)v_i=0$ for each $i=1,2,3$.

	By row reducing each equation, it can be shown that we have $v_1=\begin{pmatrix}0\\1\\0\end{pmatrix}$,
	$v_2=\begin{pmatrix}3\\-5\\1\end{pmatrix}$ and $v_3=\begin{pmatrix}0\\-1\\2\end{pmatrix}$.

	The eigenvectors of $A$ are clearly linearly independent, thus by \ref{thrm:1}, $A$ is diagonalisable, and by \ref{cor:2},
	we can write $A$ as $PDP^{-1}$ where $P=\begin{pmatrix}0&3&0\\1&-5&-1\\0&1&2\end{pmatrix}$ and $D=\begin{pmatrix}1&&\\&2&\\&&-1\end{pmatrix}$.
\end{example}
\newpage

\subsection{Case Study}

We will now


\section{Code Extracts}
% \label{sec:examples}
%
% Now we include a figure.
% (See Figure~\ref{fig:example}.)
% \begin{figure}[ht]
% 	\centering
% 	\includegraphics[width=0.3\textwidth]{example}
% 	\caption{An example of a figure}
% 	\label{fig:example}
% \end{figure}
%
% \paragraph{Acknowledgements} \lipsum[6]

%	\newpage
% \bibliography{refs}
\nocite{*}

\appendix

% \section{Omitted Proof in Section~\ref{sec:examples}}


\end{document}
